% Original: PDF strana 6, gornji slajd.
\slajdid{03-s011}
\begin{frame}[t,label=03-s011]{Aktivni nivoi i negacije}
\setlength{\parskip}{0pt}
% element: 03-s011:e01
Ulazi i izlazi mogu biti aktivni na logičkoj nuli. Primer su izlazi dekodera:
\par\vspace{.5mm}
% element: 03-s011:e02
{\centering% Sva četiri izlaza i priključci ostaju pojedinačno prikazani.
\begin{tikzpicture}[x=1mm,y=.85mm,font=\fontsize{9}{11}\selectfont]
\foreach \k/\x in {0/0,1/42,2/84}{
 \begin{scope}[shift={(\x,0)}]
 \draw[DEvod] (0,0) rectangle (24,19.8);
 \node at (12,17.6) {Dekoder};\node at (12,14.3) {2/4};
 \foreach \n/\y in {3/11.3,2/8.0,1/4.7,0/1.4}{
  \ifnum\k=1
   \node[anchor=east,inner sep=1pt] at (24,\y) {Y\n};
  \else
   \node[anchor=east,inner sep=1pt] at (24,\y) {$\overline{\mathrm{Y\n}}$};
  \fi
  \ifnum\k=0
   \draw[DEvod] (24,\y)--(27.5,\y);
  \else
   \draw[DEvod,fill=white] (24.7,\y) circle(.7);
   \draw[DEvod] (25.4,\y)--(27.5,\y);
  \fi
 }
 \node[anchor=west,inner sep=1pt] at (0,5) {CS};
 \draw[DEvod] (-3,5)--(0,5);
 \foreach \n/\xx in {1/7,0/14}{
  \node[anchor=south,inner sep=1pt] at (\xx,0) {S\n};
  \draw[DEvod] (\xx,0)--(\xx,-2.7);
 }
 \end{scope}
}
\node at (35,10) {$=$};\node[text=DEisticanje] at (77,10) {$\ne$};
\end{tikzpicture}
\par}
\par\vspace{.5mm}
Prva dva zapisa su ekvivalentna; najčešće se koristi drugi. Treći uvodi
\textcolor{DEisticanje}{\textbf{dvostruku negaciju}}, pa su izlazi ponovo aktivni na logičkoj jedinici.
\par\vspace{1mm}
% element: 03-s011:e03
{\centering
\Oznaka{Više selekcionih signala: unutrašnji signal je \(CS=CS_1\cdot\overline{CS_2}\).}
\par\vspace{.5mm}
% Prva dva CS prikaza su ekvivalentna; treći sadrži i negiran CS1.
\begin{tikzpicture}[x=1mm,y=.85mm,font=\fontsize{9}{11}\selectfont]
\foreach \k/\x in {0/0,1/42,2/84}{
 \begin{scope}[shift={(\x,0)}]
 \draw[DEvod] (0,0) rectangle (24,19.8);
 \node at (12,17.6) {Dekoder};\node at (12,14.3) {2/4};
 \foreach \n/\y in {3/11.3,2/8.0,1/4.7,0/1.4}{
  \node[anchor=east,inner sep=1pt] at (24,\y) {Y\n};
  \draw[DEvod] (24,\y)--(27.5,\y);
 }
 \draw[DEvod] (-3,4)--(0,4);
 \ifnum\k=2
  \node[anchor=west,inner sep=1pt] at (0,4) {$\overline{\mathrm{CS1}}$};
 \else
  \node[anchor=west,inner sep=1pt] at (0,4) {CS1};
 \fi
 \ifnum\k=0
  \node[anchor=west,inner sep=1pt] at (0,8) {$\overline{\mathrm{CS2}}$};
  \draw[DEvod] (-3,8)--(0,8);
 \else
  \draw[DEvod,fill=white] (-.7,8) circle(.7);
  \draw[DEvod] (-3,8)--(-1.4,8);
  \ifnum\k=1
   \node[anchor=west,inner sep=1pt] at (0,8) {CS2};
  \else
   \node[anchor=west,inner sep=1pt] at (0,8) {$\overline{\mathrm{CS2}}$};
  \fi
 \fi
 \foreach \n/\xx in {1/7,0/14}{
  \node[anchor=south,inner sep=1pt] at (\xx,0) {S\n};
  \draw[DEvod] (\xx,0)--(\xx,-2.7);
 }
 \end{scope}
}
\node at (35,10) {$=$};\node[text=DEisticanje] at (77,10) {$\ne$};
\end{tikzpicture}
\par}
\end{frame}
